% ─────────────────────────────────────────────────────────────────────────────
% The Harbor Volume Workbook — Part I: The Single-Writer Kernel (solutions)
%
% Full worked answers to the Part I questions. Check items get a tight, complete
% answer; Trace items get a step-by-step derivation; Open items get a genuine
% design discussion that takes a position, names the trade, and bounds the claim.
% Numbering matches single-writer-kernel-questions.tex (Q1..Q30).
%
% Expects \wbans{<label>}{<body>} from the workbook preamble; a fallback is below.
% ─────────────────────────────────────────────────────────────────────────────

\providecommand{\wbans}[2]{%
  \par\vspace{6pt}\noindent
  {\scshape\bfseries\textcolor{hhteal}{#1.}}\enspace #2\par\vspace{2pt}%
}

\section{The Single-Writer Kernel --- Solutions}
\label{wbsol:swk}

% ─── The seven organs ─────────────────────────────────────────────────────────
\subsection*{The seven organs}

\wbans{Q1}{Six of the seven organs survive an in-memory database; one is gutted.
Resource (claims/locks), communication (bus, markers, the two delegation objects),
obligation/enforcement (the monitor and commitments), identity, and self-attestation
all operate on live state and are meaningful with no disk: mutual exclusion,
loop-detection, deontic checks, and integrity self-checks are all about the
\emph{current} configuration, not its persistence. The \textbf{substrate organ}
becomes vacuous --- its entire job is durability and serialized mutation, and
durability with no disk is a contradiction. The \textbf{continuity organ} is the
subtler casualty: memory, recovery, and the audit log all assume state outlives the
process. ``Recovery passes notes'' presupposes the notes are still there after the
agent dies; with no disk there is nothing to pass. So: substrate vacuous outright,
continuity vacuous in spirit (it degrades to within-process bookkeeping that the next
crash erases), the other five intact.}

\wbans{Q2}{Example, the resource organ. Its home table is the claims/registry table,
keyed by the resource key (port number, file-region descriptor, or named-lock name)
with a holder column and an \texttt{expires\_at} column. The invariant it owes the
coordination layer above is \textbf{at most one live holder per key at any instant}
(I2, atomic mutual exclusion): the layer above may treat ``who holds this key'' as a
single-valued ground truth and never has to reconcile two simultaneous holders. The
table enforces this with a uniqueness/primary-key constraint on the key column, so the
invariant is a property of the schema, not of careful application code.}

\wbans{Q3}{The seven cut is defensible but not forced; it is chosen so that each
organ owns exactly one contract the layers above name. A plausible \emph{five-organ}
cut folds identity into continuity (both are ``who is this actor over time'') and
self-attestation into substrate (``the floor reports its own health''). That makes
the picture smaller and emphasizes that durability and integrity-reporting are the
same trust root. What it hides: identity is the precondition that gates the
\emph{soundness} of every enforcement rule (I12 gates I7--I9), so burying it inside
continuity obscures the single most load-bearing dependency in the threat model. A
\emph{nine-organ} cut would split the resource organ into ports / file-regions /
named-locks and the communication organ into bus / markers; this makes each table's
schema crisper but multiplies near-identical exclusion arguments and loses the
insight that all three claim granularities reduce to one exclusion primitive. The
seven-organ cut wins because it keeps one mechanism per organ while still surfacing
the identity-gates-enforcement dependency as its own thing. The deeper lesson: the
``right'' decomposition is the one whose boundaries coincide with the contracts the
next layer actually relies on, not the one with the fewest boxes.}

% ─── The substrate organ ──────────────────────────────────────────────────────
\subsection*{The substrate organ}

\wbans{Q4}{Take module A owning a \texttt{sessions} table and module B that, at a
later version, wants to read \texttt{sessions.actor\_label}. Suppose A renames
\texttt{actor} to \texttt{actor\_label} via a lazy ``add column if absent, copy,
stop reading the old one'' migration that runs when A loads. Because schema-by-union
imposes no load order, B can load and run its first read \emph{before} A has loaded
and performed the backfill: B sees a table with \texttt{actor} populated and
\texttt{actor\_label} absent or null, and either crashes or silently reads nulls. The
unsafety is precisely that a rename is a two-step (add-new, backfill, drop-old)
operation whose steps are not atomic across modules, and ``any module, any order''
gives no barrier guaranteeing the backfill precedes the first foreign read. A pure
\texttt{CREATE TABLE IF NOT EXISTS} is safe under any order because it is idempotent
and touches only its own columns; a rename is not, because it has a happens-before
obligation the union model cannot express.}

\wbans{Q5}{Walk the request \texttt{acquire(key)} from the CLI:
(1) the CLI process serializes the request and sends it over a Unix domain socket ---
no database contact yet;
(2) the daemon's single request queue receives it --- \emph{routing} point 1: the
daemon, not the OS lock, is what serializes this against other in-flight requests;
(3) the daemon, holding its one read--write connection, issues the
\texttt{INSERT}; the SQLite file lock is taken for the duration of that statement ---
this is the \emph{file-lock} point, and it is the backstop, not the primary
serializer;
(4) the constraint decides win/lose and the daemon returns the result over the socket.
So there are two routing points (socket send, daemon queue) where correctness rests on
\emph{where the write goes}, and exactly one file-lock point (the insert) where the OS
lock would still serialize even a rogue second writer. The paper's claim is that under
the discipline the file lock never actually arbitrates between two daemon writes ---
it only matters as the safety net if a second writer ever appears.}

\wbans{Q6}{Take a position: idempotent self-init can be extended a long way without a
version table, but not all the way --- down-migrations force one. The extension that
works: make each module declare not just \texttt{CREATE TABLE IF NOT EXISTS} but a
\emph{set of idempotent, commutative, monotone} schema assertions --- ``column X
exists with type T'', ``index Y exists'', ``column Z backfilled from W where null''.
Each assertion is safe in any order \emph{because} it is idempotent and only adds; the
union of monotone additions is order-independent (it is a join in a lattice of
schemas). Renames become ``add new column + backfill assertion + a read-shim that
prefers new, falls back to old'', never a destructive drop, which preserves
any-order safety at the cost of leaving dead columns. The wall you hit: a true
\emph{down}-migration (undo a column, shrink a type) is not monotone --- it removes
information --- so it cannot be an order-independent assertion, and you cannot know
whether a down-migration has already run without recording that it did. That record
is, minimally, a version/applied-migrations table. Conclusion: ``any module, any
order'' survives all forward, additive, backfill-only evolution; it cannot survive
destructive rollback, and the smallest thing that buys rollback is exactly the
sequencer the design set out to avoid. The honest design ships the additive subset and
declares destructive migration an explicit, operator-gated, sequenced operation ---
not part of the auto-init path.}

% ─── Durability by fault class ────────────────────────────────────────────────
\subsection*{Durability by fault class}

\wbans{Q7}{(a) Forced kill of the daemon process --- \textbf{I1a-survivable}: the
committed pages are in the OS page cache, which outlives the process; restart replays
the WAL. (b) Clean OS reboot --- \textbf{I1a-survivable}: a clean shutdown runs the
shutdown checkpoint (or at minimum flushes the cache), so committed data reaches
stable storage. (c) Yanking the power cord --- \textbf{I1b-exposed}: the page cache is
lost before the next checkpoint syncs the WAL, so the last committed-but-uncheckpointed
writes can roll back. (d) Unhandled exception in a request handler --- \textbf{I1a},
and in fact weaker than a kill: the process may not even die; the transaction either
committed (durable in cache) or never committed (nothing to lose). The discriminating
question for every case is the same: \emph{did the page cache survive?} Process death
and clean reboot preserve it; power loss does not.}

\wbans{Q8}{Timeline: at $t_0$ the claim commits --- its pages are in the WAL in the OS
page cache, not yet fsync'd. The vulnerable window opens at $t_0$ and closes at the
moment the WAL is durably synced, which under \texttt{synchronous=NORMAL} is the next
checkpoint at $t_0+\delta$. A power loss anywhere in the open interval $(t_0,\,
t_0+\delta)$ reverts the claim; a power loss after the checkpoint does not. Bounding
$\delta$: SQLite's \texttt{wal\_autocheckpoint} is a \emph{page-count} threshold, not
a timer --- a checkpoint fires when the WAL reaches $N$ pages (the reference config uses
$N=200$). So $\delta$ is bounded not by wall-clock but by write \emph{volume}: under a
write rate of $\rho$ pages/second the worst-case window is $\delta \le N/\rho$, and
under low write rates $\delta$ can be \emph{unbounded in time} --- a single claim on an
otherwise idle system may sit unsynced indefinitely because nothing pushes the WAL to
200 pages. This is the sharp edge: page-count checkpointing bounds the window in pages,
not in seconds, so a money-bearing write on a quiet system is exactly the case where
I1b bites hardest. The remedy (a commit-path checkpoint or full sync for that write)
is the subject of Q9.}

\wbans{Q9}{Design a \emph{per-write-path} selector rather than a global pragma. The
cleanest interface attaches a durability class to the write at the API boundary, not
to the connection:
\texttt{write(payload, durability=\{fast \textbar\ durable\})}, defaulting to
\texttt{fast}. \texttt{fast} keeps \texttt{synchronous=NORMAL} (the current behavior).
\texttt{durable} wraps the write so that immediately after the commit the daemon issues
either a \texttt{PRAGMA wal\_checkpoint(FULL)} or runs that single transaction under a
temporarily elevated \texttt{synchronous=FULL} --- forcing the WAL to stable storage
before returning success. The crucial design rule that avoids re-introducing the
conflation: the durability class must be a property of the \emph{call}, surfaced in the
return contract, so that ``success'' means different things for the two classes and the
caller chose knowingly. A settlement-ledger write opts into \texttt{durable} and pays
the fsync latency; a claim write stays \texttt{fast}. Two traps to avoid: (1) do not
make \texttt{durable} a connection-level mode that silently slows unrelated writes
sharing the connection --- it must be per-statement; (2) do not let \texttt{durable}
become the default ``to be safe'', which would erase the throughput rationale for the
single-writer design. The interface succeeds when the cost is visible at the exact
call site that needs the guarantee and nowhere else.}

% ─── The resource organ ───────────────────────────────────────────────────────
\subsection*{The resource organ}

\wbans{Q10}{Releasing a lock you do not hold returns success because release is defined
to be \textbf{idempotent} (I3): its postcondition is ``the key is not held by you,''
which already holds, so the operation is a no-op that succeeds. This is the same
property as idempotent re-claim --- the operation asserts a target state rather than a
delta. If release raised on a non-held key, two real things break. First, crash
recovery: an agent that crashed mid-release and retries on restart would get an error
on the second attempt and could not cleanly converge. Second, double-release races: two
cleanup paths (the agent's own \texttt{finally} block and the TTL sweep) legitimately
race to release the same expired lock, and exactly one would ``win'' while the other
errored --- turning a benign convergence into a spurious failure the operator must
triage. Idempotent release makes ``ensure this is released'' a safe, repeatable
command, which is what every caller actually wants.}

\wbans{Q11}{Both agents issue \texttt{acquire(K)} on fresh key $K$. Both requests reach
the single daemon and are serialized by its request queue, so the two inserts are
attempted one after the other, never truly simultaneously at the storage layer.
Step the winner (say A): \texttt{INSERT row (key=K, holder=A, expires\_at=...)}
succeeds --- $K$ had no row --- and A is returned \texttt{GRANTED(holder=A)}.
Step the loser (B), arriving an instant later:
\texttt{INSERT row (key=K, holder=B, ...)} violates the uniqueness/primary-key
constraint on \texttt{key} and is rejected \emph{atomically}; there is no instant at
which two rows for $K$ exist. Control falls to the catch path, which runs
\texttt{existing = SELECT row WHERE key=K}; this read returns A's row. The branch
\texttt{existing.holder == requester} is false (holder is A, requester is B) and the
expiry branch is false (A's lock is fresh), so B is returned
\texttt{DENIED(holder=existing.holder)} --- and \emph{that} return value is the exact
point at which the loser learns the holder's identity: it is read back from the very
row whose constraint rejected B, in the catch path, after the failed insert. B never
gets a bare error; it gets ``held by A.''}

\wbans{Q12}{Position: a FIFO wait-list of TTL'd reservations is \emph{almost} enough,
and the residual gap is exactly what forces (or avoids) a scheduler. Add a second
table, \texttt{waiters(key, requester, enqueued\_at, reservation\_ttl)}. On a denied
acquire, instead of bouncing the loser, enqueue it. On release (or TTL sweep) of the
held key, the daemon grants the key to the head of the FIFO queue for that key. This
gives bounded-wait \emph{provided} every grant is consumed --- and that proviso is the
catch. A reserved waiter that never comes back to claim its grant must time out
(\texttt{reservation\_ttl}), or the head-of-line blocks everyone behind it. Sweeping
expired reservations is the same lazy-on-read mechanism the claim table already uses,
so \emph{no background scheduler is required for fairness per se}: fairness is achieved
by ordering (FIFO) plus reservation expiry, both evaluable at request time. Where a
scheduler \emph{does} creep in: if you want \emph{priority} (not just FIFO),
\emph{anti-starvation across keys} (a slow agent waiting on many hot keys), or
\emph{proactive} grant notification (push the grant rather than wait for the waiter to
poll), you now need a component that wakes up on its own and makes allocation decisions
--- which is a scheduler, and a scheduler that decides ``who runs next'' is a
coordination-layer concern, blurring the substrate boundary the design protects. So:
plain FIFO bounded-wait with reservation TTL stays in the substrate (it is still pure
exclusion, just ordered); anything that ranks or pushes belongs one layer up. The clean
answer ships the FIFO-with-TTL primitive and exports the ordering, letting the
coordination layer build priority on top.}

% ─── The communication organ ──────────────────────────────────────────────────
\subsection*{The communication organ}

\wbans{Q13}{Read-time decay is more honest because the value a reader observes is the
value as of the read, not as of the last background tick. With tick-only decay, a
marker's stored magnitude is stale between ticks: if a tick fires every $T$ seconds and
a reader queries at $T-\epsilon$, it sees a magnitude that should have decayed for
nearly a full interval but has not. Concrete failure: a coordination marker (``agent is
actively working this file'') with a half-life shorter than the tick interval. The
depositing agent dies; no further deposits arrive; but until the next tick the stored
value is unchanged, so a reader concludes the file is hot and backs off --- coordinating
around a ghost. Read-time decay computes $w\cdot r^{\Delta t}$ with $\Delta t$ measured
from the deposit to \emph{now}, so the reader sees the correctly evaporated value
regardless of when the last tick ran; the background tick then exists only to prune
sub-threshold rows, not to keep values honest. The honesty property is: no reader ever
sees a signal fresher than physics allows.}

\wbans{Q14}{Under the silent-dedup resolution, the operator should see \emph{exactly one}
logical speech act, not two, and \emph{no} anomaly alarm. The bus redelivered the
envelope (at-least-once, working as designed), and the receiving protocol machinery
recognizes the redelivery as benign transport non-determinism and suppresses it
quietly; only a genuine protocol-state-machine violation (e.g.\ a reply to a
conversation that was already closed) would surface loudly. Trace: envelope arrives,
its effect is applied, the conversation advances; the duplicate arrives, the receiver
must determine ``have I already applied this?'' Today, without an idempotency key, that
determination is heuristic --- it leans on conversation-id plus in-reply-to plus content
equality, which can misfire (two legitimately identical acks look like a dup; a dup with
a re-stamped timestamp looks new). An envelope \textbf{idempotency key} (a stable id
minted once at first send and preserved across redeliveries) changes this from a
heuristic to an exact test: ``apply effect only if this key has not been applied,''
recorded in a small dedup table. The operator-visible behavior is identical in the happy
path; the difference is that with the key, dedup is provably correct and the rare
false-merge / false-split cases vanish.}

\wbans{Q15}{Frame it as measurement, not a guessed constant, because the ``right'' decay
is the one matched to the empirical lifetime of the coordination signal it carries, and
that lifetime differs per marker kind. Method: for each marker kind, instrument the
system to log (deposit time, the time the signal stopped being actionable). The signal
``stops being actionable'' when the situation it encodes ends --- e.g.\ for ``working
this file,'' when the agent releases the claim or dies. Collect the distribution of
these lifetimes $L$. Set the per-kind half-life so that the marker has decayed to its
prune threshold at roughly the upper quantile (say $p_{90}$) of $L$: this keeps a signal
alive for as long as 90\% of real situations last (avoiding premature evaporation) and
evaporates it shortly after (avoiding rot). Concretely, if $r^{\Delta t}$ should reach
the prune threshold $\theta/w$ at $\Delta t = L_{p90}$, solve $r = (\theta/w)^{1/L_{p90}}$
per kind. Re-measure periodically because the workload (agent count, task granularity)
shifts $L$. The deliverable is not a number but a calibration procedure plus the two
failure signals to watch (operators backing off ghosts = too slow; signals vanishing
before peers read them = too fast), which lets the rate self-correct against observed
behavior rather than a guess.}

% ─── The obligation & enforcement organ ───────────────────────────────────────
\subsection*{The obligation \& enforcement organ}

\wbans{Q16}{Walk the rules. \emph{Lock-owner validity} and \emph{capability-escalation}
checks \emph{could} move pre-commit: both are predicates on the proposed write that can
be evaluated by a before-insert trigger before the row lands (``does this releaser own
this lock?'', ``does this actor hold this capability?''), so in principle they are
regimentable --- their soundness is bounded only by I12 (a forged identity defeats the
predicate, but that is an identity gap, not a timing gap). \emph{Note monotonicity}
could also move pre-commit: ``reject any write that would lower an active session's note
count'' is a predicate on the transition, checkable at insert. The rule that genuinely
\emph{cannot} move is \textbf{heartbeat freshness}, and the reason is fundamental, not
incidental: heartbeat freshness is a \emph{failure detector} --- it concludes ``agent is
dead'' from the \emph{absence} of an event within a timeout. Chandra and Toueg show that
in an asynchronous system you cannot build a perfect failure detector: you cannot
distinguish a crashed agent from a slow one synchronously, so the judgement is
inherently retrospective (the timeout must elapse) and probabilistic. There is no insert
to reject pre-commit, because the triggering condition is a \emph{non-event} observed
after a delay. So the partition is: predicate-on-a-proposed-write rules are
regimentable; absence-of-event / timeout rules are not, and heartbeat freshness is the
canonical irreducibly-post-hoc rule.}

\wbans{Q17}{Trace the Law-2 gate with an empty oracle reference. Entry:
\texttt{close(commitment, status='done', oracle\_ref='')}. First branch:
\texttt{status != 'done'}? No --- status is \texttt{'done'}, so the drop-strategy path is
skipped. Second branch: \texttt{if oracle\_ref is empty} --- true, the reference is the
empty string. The gate returns \texttt{REFUSED("done requires a verifiable oracle
reference")} immediately. That is the exact point and reason: closure as ``done'' is
refused at the empty-reference check, \emph{before} the vocabulary classification and
\emph{before} the re-verification step, because Law 2's whole purpose is that ``done''
must be \emph{witnessed} by something the daemon can re-check, and an empty reference
witnesses nothing. The agent cannot close by assertion; it must name a released claim, a
merged commit, a passing test, or a satisfied policy sub-check. The refusal is by
construction (a branch), not by inspection of free text --- which is what makes it
Goodhart-resistant.}

\wbans{Q18}{Position and theorem sketch. Define a monitor rule $\phi$ as a predicate over
system states/transitions. $\phi$ is \textbf{regimentable} iff there exists a pure
predicate $g$ on the \emph{proposed} transition $(s, a)$ such that $g(s,a)$ is decidable
\emph{before} commit and $g(s,a) \Leftrightarrow$ ($a$ would violate $\phi$). $\phi$ is
\textbf{only enforceable} iff its violation condition references information not
available until after commit --- in particular (i) the \emph{passage of time} relative
to a future event (timeouts / failure detection), or (ii) a \emph{global aggregate} over
state the single insert does not carry. Claim: a rule is regimentable iff its violation
is a local, present-time function of the proposed write; it is merely enforceable iff
its violation is a function of elapsed time or post-hoc aggregate. Proof direction one:
if violation is a local present-time predicate, install it as a before-insert trigger;
the bad row never commits (truncation before the bad state, satisfying Schneider's
condition for enforceable safety). Direction two: if violation depends on a timeout, then
by Chandra--Toueg no synchronous pre-commit predicate decides it, so no before-insert
trigger can truncate, hence the bad state is reachable and only detect-and-compensate
remains. The deontic payoff: this partitions the monitor's rule set into a regimentable
core (move these to triggers and earn ``physically unreachable'') and an
irreducibly-post-hoc remainder (heartbeat/liveness), and it tells you that enlarging the
regimented set is exactly the work of rewriting aggregate/timeout rules as local
present-time predicates where possible.}

\wbans{Q19}{Propose a new oracle: \textbf{a signed, content-addressed artifact digest}
--- e.g.\ ``this commitment is done because artifact with hash $H$ exists in the
content-addressed store and matches the declared output contract.'' Closure supplies
$H$; the daemon re-verifies by (1) confirming an object with hash $H$ exists, and (2)
confirming $H$ satisfies a \emph{structural} contract declared when the commitment
opened (e.g.\ ``a tar of files matching glob $G$, non-empty, parsing as valid JSON
against schema $S$''). Why it is Goodhart-resistant: the agent does not get to assert
``done''; it must produce a \emph{thing} whose hash the daemon independently recomputes
and whose structure the daemon independently checks against a contract the agent could
not edit (the contract, like the deadline in Law 1, is daemon-held). The agent controls
the work, not the metric: it cannot name an artifact that does not exist (hash check
fails) nor one that does not meet the structural contract (schema check fails), and it
cannot retro-edit the contract. The line it does \emph{not} cross --- and this is the
honest limit --- is \emph{semantic} adequacy: the oracle verifies ``a well-formed
artifact of the declared shape exists,'' never ``the artifact is \emph{good}.'' That is
the same boundary as the existing vocabulary (a passing test verifies the test passed,
not that the code is correct), so the extension is consistent with the kernel's honesty
boundary rather than breaching it. Adding it widens what ``done'' can mean without
re-admitting free text, because every new condition is a re-checkable structural
predicate, not a note.}

% ─── The continuity organ ─────────────────────────────────────────────────────
\subsection*{The continuity organ}

\wbans{Q20}{``Recovery passes notes'' hands the successor agent the dead agent's durable
notes --- a summary of intent and progress. ``Recovery restores work'' would hand it the
actual execution state. They are fundamentally different because a language-model agent's
real state is not in its notes; it is in the working-tree edits in flight, the open
claims, the partially-formed plan, and the model's own in-context reasoning. One thing
notes preserve: \emph{declared intent and checkpointed conclusions} --- ``I decided to
refactor module X; tests A and B pass.'' One thing notes cannot preserve: the
\emph{uncommitted working state} --- the half-applied diff, the reasoning that had not
yet been written down, the next step the agent was about to take. A successor reading the
notes must \emph{re-derive} the in-flight work from scratch, often re-doing or conflicting
with edits the dead agent had already made on disk but not described. That gap --- between
knowing what the agent meant to do and being able to continue exactly where it stopped ---
is OP-4, and it is why checkpoint is the weakest implemented link in the continuity organ.}

\wbans{Q21}{Trace: (1) the agent calls the append-note operation; (2) it is a normal
write --- it routes through the single writer, takes the file lock for the insert, and
\textbf{commits}; the note row is now durable in the WAL/log. (3) The append emits an
event onto the append-only operation log. (4) The policy monitor, a subscriber
\emph{downstream} of that log, observes the event and evaluates the note-monotonicity
rule (``the active session's note count must not regress''). (5) If this write lowered
the count, the monitor raises a violation and compensates (flag, alert, possibly
restore). The answer to ``at which point is the violation already durable'': at step
(2) --- the offending write committed before the monitor ever saw it, because the monitor
subscribes to the post-commit log. By the post-commit enforcement theorem the monitor
cannot have prevented the regression; it can only detect it and compensate within a
bounded window. The rule is real and useful (amnesia is caught and undone) but it is
detect-and-compensate, not regimentation --- exactly the correction the paper insists on.}

\wbans{Q22}{The minimal artifact is a \textbf{content-addressed execution snapshot}, and
the proposed bundle is close to right but should be made precise. Capture, at snapshot
time: (1) the \emph{working-tree diff} against the last commit (so disk edits in flight
are recoverable); (2) the \emph{open claims} (so the successor inherits the same
exclusion footprint and does not collide); (3) the \emph{commitment set} with
daemon-held deadlines (so obligations survive the death); (4) the last $N$ turns of the
agent's \emph{transcript} (so the reasoning context can be reconstituted into a fresh
context window). Content-address the bundle (hash it) so it is tamper-evident and
dedup-able. What is recoverable from this: the exact disk state, the coordination
footprint, the outstanding promises, and enough recent reasoning to resume coherently.
What is \emph{fundamentally lost}: the model's \emph{latent} state --- the
not-yet-verbalized intuition, the half-formed next token, the implicit plan that never
reached the transcript. A transcript is a lossy projection of the model's hidden state;
you can re-prompt from it but you cannot reinstate the exact internal activations. So the
honest claim is: an execution snapshot turns ``re-derive everything'' into ``resume from
a faithful but lossy checkpoint,'' which is the difference between restarting a task and
continuing it --- and it is the literal foundation of cross-machine reputation, because a
reputation economy must be able to say ``this agent delivered this work'' even across the
agent's death, which requires the work to survive the death.}

% ─── The invariants as theorems ───────────────────────────────────────────────
\subsection*{The invariants, stated as theorems}

\wbans{Q23}{``Free'' means ``without an agreement protocol,'' not ``without cost.''
Linearizability falls out of one decider committing in a total order --- no Raft, no
quorum, no round trips. What it costs is \textbf{throughput and a scalability ceiling}.
Every mutation funnels through one writer holding one file lock, so the maximum write
rate is bounded by what a single writer can commit serially; under contention, losers
busy-wait (bounded, but still latency) and the system cannot scale writes by adding
cores or machines --- the single writer is a structural bottleneck. The trade is
deliberate and defensible for a local-first developer tool (write rates are modest;
simplicity and correctness dominate), but it is a real cost: the same architectural
choice that buys linearizability for free is also the thing that caps write
scalability, and it cannot be relaxed without re-introducing the consensus problem the
design refused. ``For free'' is true at the consistency layer and paid in full at the
throughput layer.}

\wbans{Q24}{Take three operations on key $K$:
$O_1=\texttt{acquire}(K,A)$, $O_2=\texttt{acquire}(K,B)$, $O_3=\texttt{release}(K,A)$,
issued so that in real time $O_1$ completes before $O_2$ is invoked, and $O_2$ is
invoked before $O_3$ completes but they overlap. The daemon commits in the order it
received them; suppose it received $O_1$, then $O_2$, then $O_3$. The unique
linearization is $\langle O_1: \text{GRANT } A,\; O_2: \text{DENY (held by } A),\; O_3:
\text{release } A\rangle$. This is consistent with real time: $O_1$'s linearization
point (its commit) precedes $O_2$'s, which precedes $O_3$'s, and these commit points lie
within each operation's invocation--response interval. Why no second linearization works:
any ordering that puts $O_2$ before $O_1$ would have $O_2$ \emph{grant} $B$ (the key was
free), but $O_2$ actually returned DENY, so that ordering contradicts the observed
response --- a linearization must be consistent with the values operations returned.
Likewise putting $O_3$ before $O_1$ would release a key $A$ does not yet hold; release is
idempotent so it ``succeeds,'' but then $O_1$ would still grant and the history is
unchanged --- yet real time forbids $O_3$ (invoked after $O_1$ completed) preceding
$O_1$. The single total commit order, consistent with both returned values and real-time
precedence, is unique; that uniqueness \emph{is} linearizability.}

\wbans{Q25}{The differential-test suite must compare the two SQLite-binding runtimes on
exactly the observable behavior the invariants depend on. Operation sequences to run
against both: (1) \emph{contended acquire} --- $N$ concurrent \texttt{acquire}(K) and
assert exactly one GRANT and $N-1$ DENY-with-holder, in both runtimes; (2)
\emph{expire-then-reacquire} --- acquire, let TTL pass, reacquire, asserting the
delete-then-insert sequence yields one holder and no ``database busy'' divergence; (3)
\emph{idempotent re-claim/re-release} --- repeat the same claim/release and assert no
state change and identical return; (4) \emph{WAL crash-replay} --- write, simulate
process kill, reopen, assert the committed row is present (I1a) under both bindings; (5)
\emph{pragma-semantics probe} --- set \texttt{synchronous}, \texttt{busy\_timeout},
\texttt{foreign\_keys} and assert each runtime reports and honors them identically.
Observable-state assertions: holder identity per key, return code/category for each call,
post-crash row presence, and the exact pragma values read back. The suite makes I11 a
theorem (rather than a hope) when, for every sequence, the two runtimes are proven to
produce \emph{bit-identical observable results} --- same holder, same return category,
same survivor set, same pragma readback --- because then ``green in test'' and ``correct
in production'' are the same statement. The CI consequence the paper draws follows: the
pipeline must \emph{boot the deployment runtime} and run these sequences against it, not
merely compile it.}

% ─── Cross-organ atomicity ────────────────────────────────────────────────────
\subsection*{Cross-organ atomicity}

\wbans{Q26}{Concrete partial failure: a ``begin session'' that should (1) claim port
7421, (2) open session $S$, (3) record commitment $C$. Steps (1) and (2) commit, then
(3) fails (say a constraint or a transient error). Inconsistent state left behind: port
7421 is now held and session $S$ is open, but there is \emph{no} commitment $C$ binding
the work --- so the agent holds a resource and a session with no obligation governing
them, and from the operator's view a session exists that the obligation organ does not
know about. Who sees it: the \emph{agent} sees a half-initialized session (it thinks
``begin'' failed and may retry, double-claiming or orphaning); the \emph{operator} sees
a port held and a session open with no commitment, which the monitor cannot reconcile
because no rule says ``a session must have a commitment.'' A \emph{second agent} sees
port 7421 as taken and cannot use it, even though the owning session is effectively
dead-on-arrival. The damage is bounded (it is all local rows) but it is exactly the
undefined partial-failure semantics §12 names --- and Q27's single transaction erases
it.}

\wbans{Q27}{Interface change: wrap the three writes in one local transaction at the
endpoint boundary. Pseudocode: \texttt{BEGIN IMMEDIATE; claim(port); open(session);
record(commitment); COMMIT}, with any failure triggering \texttt{ROLLBACK} so the
operation is all-or-nothing. \texttt{IMMEDIATE} is required so the write lock is taken at
\texttt{BEGIN} rather than lazily on first write, which closes the busy-error window the
paper flags in the stale-holder branch. Because all three writes are local rows in one
SQLite file under one writer, this single transaction gives atomicity, consistency, and
isolation for the composite operation for free --- no compensating logic, no intermediate
visible state. Contrast with Sagas: if you instead chose a saga, you would execute the
three steps as separate committed transactions and attach a \emph{compensating action} to
each --- release-port compensates claim-port, close-session compensates open-session,
retract-commitment compensates record-commitment --- and on failure of step 3 you would
run the compensations for steps 2 and 1 in reverse. The saga is strictly more complex
\emph{and strictly weaker}: it exposes intermediate states (other agents can observe the
port held before the saga aborts), the compensations can themselves fail, and you must
reason about every interleaving. The single transaction dominates because Sagas exist
precisely for the case the kernel does \emph{not} have --- remote steps where one
transaction is impossible. Here every step is a local write, so the simpler, stronger
primitive applies. Filing this under ``research'' overstated it: it is a buildable
defect, and the build is one \texttt{BEGIN IMMEDIATE} wrapper.}

% ─── Threat model ─────────────────────────────────────────────────────────────
\subsection*{Threat model and the limits of mediation}

\wbans{Q28}{The hash chain defends against a \textbf{tamperer who is not the same user}
--- specifically, a process or party that can present or replicate the audit log across
a trust boundary the local owner does not control. Under the substrate's stated threat
model the same-user process is out of scope (it already owns the file), so on a single
machine the chain protects against nobody real. The adversary becomes real at the
\textbf{economy layer}: when an audit log is shipped across machines to another operator,
that operator cannot assume the sender's good faith, and the hash chain lets them verify
the log was not retroactively edited in transit or by the sender. So the chain is
``provisioned locally, activated across machines'': built now, load-bearing only when a
cross-operator must trust a log it did not produce. Naming the out-of-scope adversary is
the honest move --- a control that defends nobody in the current model is re-scoped, not
quietly claimed as security.}

\wbans{Q29}{Trace the forgery. The lock-owner-validity rule says ``a release of lock $L$
succeeds only if the requester's actor identity equals $L$'s recorded holder.'' Agent $X$
wants to release agent $Y$'s lock. (1) $X$ constructs a request asserting actor identity
$Y$ --- and here is the gap: identity is \emph{self-asserted} (I12 is \textsc{designed},
not built), so the session simply claims to be $Y$ and nothing cryptographic contradicts
it. (2) The release request reaches the daemon; the lock-owner-validity check compares the
\emph{claimed} requester ($Y$) against the recorded holder ($Y$) --- they match. (3) The
release is granted; $Y$'s lock is gone, released by $X$. The exact point the forgery slips
through is step (1)/(2): the validity rule keys on an identity field that the requester
populates and that no signature binds to a real principal. The rule is sound \emph{given}
an unforgeable identity and unsound without one --- which is precisely the paper's claim
that I12 gates the soundness of I7--I9. Cryptographic identity (a per-actor signing key,
with the release request signed) would make step (1) impossible: $X$ cannot produce $Y$'s
signature, so the claimed identity no longer matches a verifiable principal.}

\wbans{Q30}{Layered hardening against the same-machine adversary, smallest first, with
each layer's stopping point. (1) \textbf{OS keychain for the signing key} --- move the
private key out of the owner-readable database file into the platform keychain
(Keychain/Secret Service/DPAPI), so a same-user process that reads the DB no longer gets
the key. Stops: casual key exfiltration via file read. Does not stop: a process running
\emph{as the same user} that asks the keychain to sign on its behalf, since keychain ACLs
are typically per-user, not per-process. (2) \textbf{Per-process keychain ACLs / signing
broker} --- gate signing behind a broker that authenticates the calling process, so only
the daemon binary can use the key. Stops: arbitrary same-user processes minting
signatures. Does not stop: a process that injects into or impersonates the daemon binary.
(3) \textbf{Sealed memory / OS sandbox} (Landlock, \texttt{unveil}, Seatbelt) around the
daemon --- confine the daemon's own surface and seal the in-memory key so even a debugger
on the same uid cannot scrape it. Stops: memory scraping and ptrace of the live daemon.
Does not stop: a privileged (root) adversary or hardware attack. The honest boundary: each
layer raises the bar for a same-user attacker, and the sequence asymptotes at ``everything
short of root and hardware'' --- which is exactly where a TPM/secure-enclave dependency
would begin. The minimal \emph{useful} hardening is layers (1)+(2): key in the keychain
behind a process-authenticated broker. That closes the realistic same-user adversary
(another of your own tools reading the file) without the TPM dependency the question
rules out; sealing memory is the next increment when the threat model promotes the
same-uid-injection attacker into scope.}
